\documentclass[11pt]{article}

\usepackage[margin=1.2in]{geometry}
\usepackage{amsmath,amssymb}
\usepackage{microtype}
\usepackage{booktabs}
\usepackage{xcolor}
\usepackage{mdframed}
\usepackage[colorlinks=true,linkcolor=blue!50!black,urlcolor=blue!50!black]{hyperref}

\newmdenv[backgroundcolor=blue!4,linecolor=blue!40,linewidth=1pt,
  skipabove=8pt,skipbelow=8pt]{keyidea}

\title{The Word-Complexity Ladder:\\
       \Large A Friendly Guide}
\author{Companion notes to ``A Machine-Verified Rung Zero and a
Mahler-Value Reduction of the Automatic Rung''\thanks{Written in
collaboration with Claude (Anthropic). These notes explain the ideas;
the paper and the Lean files contain the proofs. The Collatz
conjecture remains open.}}
\date{June 12, 2026}

\begin{document}
\maketitle

\section{The diary of an orbit}

Play the Collatz game from any number: halve if even,
triple-and-add-one-then-halve if odd. Alongside the numbers
themselves, keep a \emph{diary}: write $1$ each time the number was
odd, $0$ each time it was even. The diary of $3$ reads
$1\,1\,0\,0\,0\,1\,0\,1\,0\,1\ldots$ --- after five entries the orbit
reaches the $1 \to 2 \to 1$ loop, and the diary settles into the
repeating pattern $1\,0\,1\,0\ldots$ forever. Diaries of settled
orbits always end in a repeating pattern.

Two facts from the earlier papers in this series set the stage. First,
\textbf{the diary determines the number}: no two starting numbers ever
write the same infinite diary. (Each new entry reveals one more binary
digit of where you started --- after $T$ entries you know the start
modulo $2^T$.) Second, \textbf{a number that escapes to infinity must
keep its diary at least $63.09\%$ ones forever} --- the critical
density $\log 2/\log 3$ that runs through this whole subject. Below
that share of odd steps, the halvings win and the orbit falls.

So here is a strategy for attacking the divergence half of the
conjecture, one class of diaries at a time:

\begin{keyidea}
\textbf{The ladder.} Sort all possible diaries by how complicated they
are: perfectly repetitive ones at the bottom; then diaries a simple
machine could write; then more inventive ones; with truly patternless
diaries at the top. Prove, rung by rung, that \emph{no actual number
writes a diary that simple}. Every rung is a finite, honest theorem.
The conjecture's divergence half is the (unreachable-for-now) top.
\end{keyidea}

\section{Rung zero: repeating diaries belong to settled orbits}

The floor of the ladder is now a machine-checked theorem: \emph{a
diary that eventually repeats belongs to a bounded orbit}. Why:
if the diary from day $s$ onward equals the diary from day $s+P$
onward, then the two orbit values on those days have \emph{identical
infinite diaries} --- and since diaries determine numbers, the values
are equal. The orbit literally loops, so it can never have been
heading to infinity. Every step of this argument is formalized in
Lean (no axioms, no gaps), including the bookkeeping identity that
powers everything else below.

An escaping number, then, must improvise forever --- its diary can
never settle into a repeat. The next rung up asks: could the
improvisation at least be \emph{mechanical}?

\section{Rung one: diaries written by a copying machine}

Here is the simplest kind of aperiodic diary. Take the two rules
\[
1 \;\to\; 110, \qquad 0 \;\to\; 011,
\]
start with $1$, and repeatedly rewrite every symbol by its rule:
\[
1 \;\to\; 110 \;\to\; 110\,110\,011 \;\to\; \cdots
\]
Each pass triples the length and the text stabilizes into one
infinite word $w = 110110011110110011011110110\ldots$ --- never
repeating (the paper gives a four-line proof), yet generated by a
two-line program. Words like this are called \emph{automatic}. Note
its density of ones: exactly $2/3 = 66.7\%$, comfortably above the
critical $63.09\%$. As far as the density theory is concerned, this
diary is \emph{allowed} to belong to an escaping number. Does it?

By fact one, at most one number could own this diary. Which number?
The bookkeeping identity of the earlier papers answers exactly: the
candidate owner is the value of the series
\[
x \;=\; -10 \;+\; \tfrac{5}{9}\,F\!\left(\tfrac{8}{9}\right),
\qquad\text{where}\quad
F(z) \;=\; w_0 + w_1 z + w_2 z^2 + w_3 z^3 + \cdots
\]
--- the ``diary polynomial'' of infinite degree, evaluated at the
innocent-looking fraction $8/9$. (Why $8$ and $9$? Three game steps
multiply the number by roughly $\tfrac{3 \cdot 3}{2\cdot 2\cdot 2}$
when two of them are odd --- the $8$ counts the halvings in a block
of three steps, the $9$ counts the triplings. The fraction $8/9$
\emph{is} the substitution rule, seen by arithmetic.)

One catch, and it is the heart of the matter: the series must be
summed \emph{$2$-adically} --- in the number system where a number is
small when it is divisible by a high power of $2$ (the natural habitat
of Collatz diaries, which are read off by repeated halving). The same
series also sums happily as ordinary real numbers --- to
$-6.5701\ldots$, visibly not a positive integer --- but that tells us
nothing: the real sum and the $2$-adic sum of the same series are
genuinely different numbers. It is the $2$-adic value that would be
the escaping integer.

\begin{keyidea}
\textbf{The whole rung in one sentence.} The infinite word $w$ copies
itself under the rule --- so its diary polynomial satisfies the exact
self-similarity equation
\[
F(z) \;=\; \frac{z+z^2}{1-z^3} \;+\; (1-z^2)\,F(z^3),
\]
and the only number that could own the diary is
$-10 + \frac59 F(8/9)$ summed $2$-adically. So: \emph{if $F(8/9)$,
summed $2$-adically, is not a fraction, then no whole number writes
this diary} --- one rung of the ladder hangs on the irrationality of
one exotic number.
\end{keyidea}

\section{A question with a famous shape}

``Self-similar function, evaluated at a nice point: can the value be
rational?'' is a century-old genre of question. Kurt Mahler invented
a method for it in 1929, and the modern theory (Nishioka's book is
the standard reference; Adamczewski and Faverjon's recent work its
sharpest form) handles exactly equations like the one above. Indeed,
for our $F$ the classical theory already delivers a verdict at the
\emph{real} number side: $F(8/9)$ summed as real numbers is provably
transcendental --- not just irrational, but as far from a fraction as
a number can be. The hypotheses are checked in the paper: the word
never repeats, so $F$ is a genuinely transcendental \emph{function}
(a lovely 1906 theorem of Fatou does this step), and the
self-similarity equation is of precisely Mahler's kind.

What rung one needs is the \emph{same} conclusion for the \emph{same}
function at the \emph{same} point --- but summed $2$-adically. The
ingredients of Mahler's proof are not specifically about real
numbers, and analysts have transposed his method to other number
systems before; but a published theorem in exactly the needed form
has not yet been located. That is the paper's punchline: an honest,
well-posed question handed to the right specialists.

\section{Why the lazy route provably fails}

Couldn't one just say: the series converges \emph{fast}
$2$-adically, and numbers that admit too-fast rational approximation
are irrational (Liouville's classic trick)? The paper computes
exactly how fast: cutting the series after $N$ terms approximates the
value with quality exponent
\[
\frac{3\log 2}{2\log 3} \;=\; 0.9464\ldots
\]
--- and Liouville-style arguments need quality strictly above $1$.
Worse (or better): for \emph{any} candidate diary in the allowed
band, the quality works out to (critical density)/(diary's density),
which is below $1$ precisely because escaping diaries must beat the
critical density. The same constant that creates the divergence
problem seals off the cheap proof. Crossing from $0.946$ to beyond
$1$ is exactly what Mahler's auxiliary-polynomial machinery was built
to do --- speed alone won't, structure must.

\section{The fairness check}

A recurring discipline in this series: any method claiming progress
on $3x+1$ must \emph{not} accidentally ``work'' for the $5x+1$ game,
where orbits really do seem to escape (start at $7$ and watch). The
ladder passes. For $5x+1$ the same construction evaluates the same
diary polynomial at $8/25$ instead of $8/9$, and the Mahler argument
would prove: no whole number writes \emph{this exact mechanical
diary} under $5x+1$ either. That is consistent with real escape ---
the actual $5x+1$ escapees write messy, patternless diaries (about
$49\%$ ones), nothing like a copying machine's output. The method
excludes specific orderly diaries; it cannot fake a proof that
everything settles. That is what soundness looks like here.

\section{What this does and does not achieve}

Honestly stated: even a full success at rung one would exclude only
the most orderly aperiodic diaries --- a countable, measure-zero
family --- and the conjecture would remain open above it. The earlier
no-go theorem in this series proves the top of the ladder cannot be
reached by any argument with bounded memory; the ladder is a way of
making clean partial progress, not a secret staircase to the summit.

What the paper does deliver: the floor of the ladder machine-verified;
the second rung converted --- by exact, mostly machine-checked
arithmetic --- into a single, classical-shaped question about one
$2$-adic number; a proof that the question genuinely cannot be
settled by elementary approximation, with the obstruction quantified
($0.946 < 1$) and traced back to the problem's own critical constant;
and a fairness certificate that the method cannot overclaim. Each
piece is checkable: the formal parts by running Lean, the numerical
parts by a one-minute script.

\section{Summary in four sentences}

Every Collatz orbit writes a binary diary of its parities, the diary
determines the orbit, and escaping orbits must keep their diaries
above $63.09\%$ ones forever --- so excluding diaries class by class,
ordered by complexity, is a ladder of honest theorems toward the
divergence half of the conjecture. Rung zero is machine-verified:
repeating diaries belong to settled orbits only. Rung one --- diaries
written by the simplest copying machines --- reduces exactly to
whether one self-similar function's value at $8/9$, summed
$2$-adically, is irrational: the real-number version is already
transcendental by Mahler's classical method, the cheap approximation
route provably falls short by the critical-density margin, and the
$2$-adic case is a well-posed open question for transcendence theory.
The conjecture remains open --- but one of its rungs now has the
shape of a century-old, well-understood kind of problem.

\end{document}
